\documentclass[12pt]{article}
\usepackage[T1]{fontenc}
\usepackage[utf8]{inputenc}
\usepackage{lmodern}
\usepackage{textcomp}
\usepackage{enumitem}
\usepackage{geometry}
\geometry{margin=1in}
\begin{document}
\begin{center}
Answer Key - Geography - Undergraduate Level Assessment 7 \
\small (Answers are ordered exactly as in the question paper.)
\end{center}

\section*{Multiple Choice}
\begin{enumerate}[label=\arabic*.]
\item \textbf{Answer:} B
\\ \textit{Options:} A) 19\%; B) 34\%; C) 49\%; D) 64\%
\\ \textit{Note:} The correct answer is 34\% because, according to the 2011 Census of India, approximately one-third of the Indian population resided in urban areas, and this percentage has remained consistent in subsequent estimates leading up to 2017.
\item \textbf{Answer:} A
\\ \textit{Options:} A) minimum wage legislation; B) land reform; C) progressive taxation; D) increased access to education
\\ \textit{Note:} Minimum wage legislation is not an element of the redistribution-with-growth policy approach because this approach primarily focuses on strategies that promote economic growth while simultaneously addressing income inequality, such as social investment and progressive taxation, rather than directly imposing wage controls.
\item \textbf{Answer:} B
\\ \textit{Options:} A) 41\%; B) 56\%; C) 71\%; D) 86\%
\\ \textit{Note:} The correct answer of 56\% reflects the significant value Italians place on free media as a fundamental aspect of democratic society, indicating a strong desire for the protection of journalistic independence from government interference.
\item \textbf{Answer:} C
\\ \textit{Options:} A) 18\%; B) 38\%; C) 58\%; D) 78\%
\\ \textit{Note:} The correct answer of 58\% reflects the significant growth in internet accessibility and usage in Brazil by 2017, indicating that a majority of the population had integrated internet into their daily lives.
\item \textbf{Answer:} B
\\ \textit{Options:} A) 55,000; B) 550,000; C) 5,500,000; D) 55,000,000
\\ \textit{Note:} The correct answer of 550,000 reflects the estimated number of homeless individuals in the United States as reported by the U.S. Department of Housing and Urban Development in their annual point-in-time count conducted in January 2016.
\end{enumerate}

\section*{True / False}
\begin{enumerate}[label=\arabic*.]
\setcounter{enumi}{5}
\item \textbf{Answer:} False
\\ \textit{Options:} A) True; B) False
\\ \textit{Note:} The correct answer is: 12
\item \textbf{Answer:} False
\\ \textit{Options:} A) True; B) False
\\ \textit{Note:} The correct answer is: 6 th century BC
\item \textbf{Answer:} True
\\ \textit{Options:} A) True; B) False
\\ \textit{Note:} According to the passage: 3600
\item \textbf{Answer:} False
\\ \textit{Options:} A) True; B) False
\\ \textit{Note:} The correct answer is: 2 nd
\item \textbf{Answer:} False
\\ \textit{Options:} A) True; B) False
\\ \textit{Note:} The correct answer is: the second half of the 6 th century
\end{enumerate}

\section*{Long Form Response}
\begin{enumerate}[label=\arabic*.]
\setcounter{enumi}{10}
\item \textbf{Answer:} the A$_{1}$
\\ \textit{Note:} Expected answer: the A$_{1}$
\item \textbf{Answer:} December 13, 2011
\\ \textit{Note:} Expected answer: December 13, 2011
\end{enumerate}

\vfill
\noindent\textit{End of Answer Key}
\end{document}